% Einheit 2 von 4 – Prozentwert berechnen (Katalog Einheit 3)
\einheitenkopf[e2][2 Prozentwert]{Einheit 2 von 4 · Prozentwert berechnen}
\zweigzeile{Hier lernst du auszurechnen, wie viel ein Anteil in Prozent von einer Größe ist · neu in diesem Jahr · P10 · baut auf: Streifen einteilen (Nr.~13), Bruch als Dezimalzahl (Nr.~5), durch 100 teilen (Nr.~7), Dreisatz (Nr.~8), Bruchteil einer Größe (Nr.~9)}
\setcounter{aufgabe}{23}

\verfahren{Teil am Streifen}
\begin{aufgabe}{Ich kann am Streifen ausrechnen, wie viel ein einfacher Anteil in Prozent von einer Größe ist.}
Rechne am Streifen. Der gesuchte Teil heißt Prozentwert.

\swz{$50\,\%$ von $20\,€$}{\streifen[2]{50}{$0\,€$}{$20\,€$}}
\swz{$50\,\%$ von $60\,€$}{\streifen[2]{50}{$0\,€$}{$60\,€$}}
\swz{$50\,\%$ von $140\,€$}{\streifen[2]{50}{$0\,€$}{$140\,€$}}
\swz{$50\,\%$ von $300\,€$}{\streifen[2]{50}{$0\,€$}{$300\,€$}}
\swfrage{Was bleibt gleich, was ändert sich?}
\end{aufgabe}

\begin{aufgabe}{Ich kann am Streifen ausrechnen, wie viel ein einfacher Anteil in Prozent von einer Größe ist – weiter.}
Rechne am Streifen.

\swz{$25\,\%$ von $60\,$kg}{\streifen[4]{25}{$0\,$kg}{$60\,$kg}}
\swz{$10\,\%$ von $70\,$m}{\streifen{10}{$0\,$m}{$70\,$m}}
\swz{$25\,\%$ von $30\,$kg}{\streifen[4]{25}{$0\,$kg}{$30\,$kg}}
\end{aufgabe}

\verfahren{Teil über 1 \% und über 10 \%}
\begin{aufgabe}{Ich kann über 1 \% ausrechnen, wie viel ein Anteil in Prozent von einer Größe ist.}
Rechne mit dem Dreisatz über 1 \%.
\rechenplatz{3}

\swz[1]{$7\,\%$ von $300\,€$}{\begin{dreisatz}{Prozent}{Euro}
\dsz{100\,\%}{300\,€}
\dsp{:100}{:100}
\dsz{1\,\%}{\dsleer}
\dsp{\cdot 7}{\cdot 7}
\dsz{7\,\%}{\dsleer}
\end{dreisatz}}
\swz[1]{$6\,\%$ von $50\,$kg}{\begin{dreisatz}{Prozent}{Kilogramm}
\dsz{100\,\%}{50\,\text{kg}}
\dsp{\phantom{:0}}{\phantom{:0}}
\dsz{1\,\%}{\dsleer}
\dsp{\phantom{:0}}{\phantom{:0}}
\dsz{6\,\%}{\dsleer}
\end{dreisatz}}
\end{aufgabe}

\begin{aufgabe}{Ich kann über 10 \% ausrechnen, wie viel ein Anteil in Prozent von einer Größe ist.}
Rechne über 10 \%.

\swz{$30\,\%$ von $90\,€$}{\streifen{30}{$0\,€$}{$90\,€$}}
\swz{$15\,\%$ von $60\,€$}{\streifen[20]{15}{$0\,€$}{$60\,€$}}
\swz{$5\,\%$ von $140\,€$}{\streifen[20]{5}{$0\,€$}{$140\,€$}}
\end{aufgabe}

\verfahren{Teil mit der Dezimalzahl}
\begin{aufgabe}{Ich kann einen Anteil in Prozent als Dezimalzahl mit einer Größe malnehmen.}
Schreibe den Prozentsatz als Dezimalzahl und rechne.

\swz{$40\,\%$ von $30\,€$}{$40\,\% = \rule{12mm}{0.4pt}$ \qquad $\rule{12mm}{0.4pt} \cdot 30\,€ = \rule{12mm}{0.4pt}$}
\swz{$8\,\%$ von $250\,€$}{$8\,\% = \rule{12mm}{0.4pt}$ \qquad $\rule{12mm}{0.4pt} \cdot 250\,€ = \rule{12mm}{0.4pt}$}
\swz{$20\,\%$ von $4{,}50\,€$}{$20\,\% = \rule{12mm}{0.4pt}$ \qquad $\rule{12mm}{0.4pt} \cdot 4{,}50\,€ = \rule{12mm}{0.4pt}$}
\end{aufgabe}

\verfahren{Teil im Sachtext: sparen und bezahlen}
\begin{aufgabe}{Ich kann ausrechnen, wie viel man spart und wie viel man noch bezahlt.}
Lies genau, was gefragt ist, und rechne.

\swz[3]{Ein Rucksack kostet $60\,€$. Es gibt $25\,\%$ Rabatt. Wie viel Euro spart man?}{\streifen[4]{25}{$0\,€$}{$60\,€$}}
\swz[3]{Wie viel muss man für den Rucksack noch bezahlen?}{\streifen[4]{25}{$0\,€$}{$60\,€$}}
\swz[3]{Ein Fahrrad kostet $350\,€$. Es gibt $10\,\%$ Rabatt. Wie viel kostet es jetzt?}{\streifen{10}{$0\,€$}{$350\,€$}}
\begin{teile}
\teil Ein Tablet kostet $380\,€$. Im Angebot gibt es $15\,\%$ Rabatt. Wie viel Euro spart man? Kreuze an. (P10 2021 OS)\\
\kreuz{$5{,}70\,€$}\quad \kreuz{$15\,€$}\quad \kreuz{$57\,€$}\quad \kreuz{$323\,€$}
\end{teile}
\end{aufgabe}

\begin{aufgabe}{Ich kann ausrechnen, wie viel mehr als 100 \% von einer Größe sind (je nach Buch erst in Klasse 8).}
Rechne.

\swz{$150\,\%$ von $40\,€$}{\streifen[2]{100}{$0\,€$}{$40\,€$}}
\swz[3]{Eine Sonnenblume war im Juni $80\,$cm hoch. Im August ist sie $130\,\%$ so hoch. Wie hoch ist sie jetzt?}{\streifen{100}{$0\,$cm}{$80\,$cm}}
\end{aufgabe}

\begin{aufgabe}{Ich finde den Fehler, wenn jemand einen Anteil in Prozent von einer Größe ausrechnet.}
Finde den Fehler, benenne ihn und rechne richtig.

\swz[3]{$40\,\%$ von $60\,€$. Jonas rechnet so:}{\rechnung{0{,}04 \cdot 60\,€ &= 2{,}40\,€}}
\swz[3]{Eine Hose kostet $45\,€$. Es gibt $20\,\%$ Rabatt. Wie viel Euro spart man? Emma rechnet so:}{\rechnung{0{,}8 \cdot 45\,€ &= 36\,€}}
\end{aufgabe}

\begin{aufgabe}{Ich kann begründen, warum man für 1 \% durch 100 teilt.}
Begründe.

\swfrage{Warum teilt man eine Größe durch 100, um $1\,\%$ davon zu bekommen?}
\swfrage{Lara sagt: „$20\,\%$ von $50\,€$ ist genauso viel wie $50\,\%$ von $20\,€$.“ Stimmt das? Begründe ohne zu rechnen.}
\end{aufgabe}

\begin{aufgabe}{Ich kann mit der Mehrwertsteuer Preise vergleichen.}
Ein Handwerker berechnet für eine Reparatur $240\,€$ ohne Mehrwertsteuer. Dazu kommen $19\,\%$ Mehrwertsteuer. Ein zweiter Handwerker verlangt $280\,€$ mit Mehrwertsteuer. Rechne und entscheide.

\swz[1]{Wie viel Euro Mehrwertsteuer kommen beim ersten Handwerker dazu?}{\begin{dreisatz}{Prozent}{Euro}
\dsz{100\,\%}{240\,€}
\dsp{\phantom{:0}}{\phantom{:0}}
\dsz{1\,\%}{\dsleer}
\dsp{\phantom{:0}}{\phantom{:0}}
\dsz{19\,\%}{\dsleer}
\end{dreisatz}}
\swz{Wie viel kostet die Reparatur beim ersten Handwerker insgesamt?}{}
\swfrage{Welcher Handwerker ist günstiger? Begründe.}
\end{aufgabe}

\uebersichtskasten{Prozentwert: erst 1 \% (Ganzes : 100), dann mal p – oder p als Dezimalzahl mal Ganzes. \\ 18 \% von 40 kg: \quad 1 \% $= 0{,}4$ kg $\rightarrow$ 18 \% $= 7{,}2$ kg \qquad $0{,}18 \cdot 40 = 7{,}2$}
